\section{What The Backward Walk Carries Forward}

The chapter has walked the tower back to the needle. It did not erase the
forward construction. It returned through it. The reader's predicates gave
way to quotient selection. Quotient selection gave way to boundary
orientation. Boundary orientation gave way to loop holonomy. Loop holonomy
gave way to charge. Charge gave way to the invariant and the finite gauge.
The finite gauge returned to the first sanctioned identity. The loop came
home.

That closure made charge readable as loop count. An open path may carry
transport, but it does not charge the ledger until it returns to its
reference. The closed traversal is what can be counted, and the integer
character of the count follows from the whole return. Charge is the lower
face of the number because it reads completed traversal before it reads
strain.

Mass then appeared as the upper face. The grammar carried null, threshold,
and response until the second difference became visible. That second
difference is strain, and strain is the cost read as mass. The chapter did
not inject mass behind the ledger. It read mass where the loop's response
bent past threshold.

The one number then showed four faces. Charge and mass were projections:
lower loop count and upper strain. Phase was a decision face, carrying how
orientation is read under the permitted baseline. Sameness was the
identification face, allowed to touch the needle only because closure had
restored one bounded identity. The faces differed by cost, but they remained
faces of one loop.

The imaginary bounce sharpened the phase face. The grammar read the
quarter-turn as an orientation cost and read the sign split through repeated
phase return. It kept the claim disciplined: the bounce is a phase reading
on the carriers already earned, not a completed construction of every later
spinor object. Flat closure reads the native negative sign. Tilted closure
can read the positive sign.

Finally the electron reached the Cooper-pair boundary. There the admissible
transport is paired. The boundary permits charge in units of \(2e\), forbids
incompatible interior strain, carries exactly one positive gauge turn in the
paired reading, and charges phase cost for coherent closure. The
superconducting examples illustrate the interface, but the grammar supplies
the obligation.

Chapter 09 receives this result: a loop-counted, strain-bearing,
phase-oriented number standing at a boundary where charge is no longer read
only as a solitary traversal. The backward walk has closed the tower around
the needle and brought the electron to the Cooper-pair ledger. The next
question is how far a finite reading can extrapolate from that boundary
without forgetting the bounds that made the count honest.
